\small\begin{longtable}{>{\raggedright\arraybackslash}p{.09\textwidth}>{\raggedright\arraybackslash}p{.39\textwidth}>{\raggedright\arraybackslash}p{.20\textwidth}>{\raggedright\arraybackslash}p{.21\textwidth}}
\caption{Seleção cronológica de títulos inéditos no inventário anterior: uma referência por obra, periódico ou site. Datas de captura são explicitadas, sem retrodatação.}\\
\toprule Ano / ID & Título e localização & Forma / categoria & Base da data \\\midrule\endfirsthead
\toprule Ano / ID & Título e localização & Forma / categoria & Base da data \\\midrule\endhead
1645 / N01 & Chronica da Companhia de Iesv na provincia de Portvgal [...] (Volume 1); PDF p.266 \cite{r4chronicacompanhiaiesvportugal} & risos (normalização do s longo); lexical & ano do exemplar no catálogo; não primeira edição presumida\\
1720 / N02 & Vocabulario Portuguez e latino (Volume 07: Letras Q-S); PDF p.353 \cite{r4vocabularioportuguezlatino} & Riso; verbete, p. impressa 345; lexicográfica & ano do exemplar no catálogo; não primeira edição presumida\\
1789 / N03 & Diccionario da lingua portugueza composto pelo padre D. Rafael Bluteau, reformado, e accrescentado por Antonio de Moraes Silva natural do Rio de Janeiro (Volume 2: L - Z); PDF p.354 \cite{r4diccionariolinguaportuguezamorais} & RISO; verbete; lexicográfica & ano do exemplar no catálogo; não primeira edição presumida\\
1798 / N04 & Viola de Lereno : collecção das suas cantigas, offerecidas aos seus amigos. (Volume I); PDF p.259 \cite{r4violadelereno} & Mostra em riso moderado; lexical & ano do exemplar no catálogo; não primeira edição presumida\\
1813 / N05 & O Patriota, jornal litterario, político, mercantil (v. 1, t. 2, n. 4-6); PDF p.25 \cite{r4opatriota} & De riso menos digno; lexical & ano do exemplar no catálogo; não primeira edição presumida\\
1822 / N06 & Gazeta Pernambucana; PDF p.26 \cite{r4gazetapernambucana} & isso merece risada; lexical & ano do exemplar no catálogo; não primeira edição presumida\\
1836 / N07 & Nitheroy, revista brasiliense, sciencias, letras e artes (f.1); PDF p.140 \cite{r4nitheroy} & riso de alegria; lexical & ano do exemplar no catálogo; não primeira edição presumida\\
1837 / N08 & Reflexões dramaticas para uso dos canditados que se dedicam a scena; PDF p.22 \cite{r4reflexoesdramaticas} & evitando sempre o rir-se a gargalhadas; lexical & ano do exemplar no catálogo; não primeira edição presumida\\
1846 / N09 & Primeiros cantos : poesias; PDF p.51 \cite{r4primeiroscantos} & Da casa do guerreiro o riso, a dança; lexical & ano do exemplar no catálogo; não primeira edição presumida\\
1866 / N10 & Gennesco: vida academica; PDF p.12 \cite{r4gennesco} & philosophia commodista dos risos; lexical & ano do exemplar no catálogo; não primeira edição presumida\\
1902 / N11 & A fallencia; PDF p.12 \cite{r4afallencia} & ás gargalhadas; lexical & ano do exemplar no catálogo; não primeira edição presumida\\
1920 / N12 & Negrinha (contos); PDF p.23 \cite{r4negrinha} & torceram-se de riso; lexical & ano do exemplar no catálogo; não primeira edição presumida\\
1980 / N13 & Poesia, 1958/1978; PDF p.35 \cite{r4poesia19581978} & O riso tem lágrimas; lexical & ano do exemplar no catálogo; não primeira edição presumida\\
2007 / N14 & Kibe Loco: Uma chance (parte 2); postagem Uma chance (parte 2); parágrafo com não contendo o riso \cite{r4kibeloco} & não contendo o riso; lexical & data exibida na postagem; captura de 01/07/2007\\
2010 / N15 & Cifra Club: como afinar meio tom abaixo; resposta Enviado: 12/jul/10 15:21; após boa! \cite{r4forumcifraclub} & kkkkkkkkkkkkk; mimética & data e hora exibidas na mensagem; captura de 15/09/2010; tópico iniciado em 2003\\
2014 / N16 & Orkut: comunidade Kiko Dança Psy-TRANCE; descrição da comunidade, após sempre agitado \cite{r4orkutcomunidades} & hehehe; mimética & data da captura, não data de redação; criação da comunidade em 2006 não data o trecho\\
2017 / N17 & Rubens Borba e Moraes: Anotações de um Bibliófilo, Cristina Antunes; PDF p.146 \cite{r4rubensborbamoraesanotacoesbibliofilo} & cujas poezias serviam de riso; lexical & ano do exemplar no catálogo; não primeira edição presumida\\
\bottomrule\end{longtable}\normalsize
